\documentclass[11pt]{article}

%==============================================================
% Morphogenetic Aspect Networks (arXiv-style preprint, draft)
% Single-column article. Self-contained: no external .sty/.bst.
%==============================================================

\usepackage[utf8]{inputenc}
\usepackage[T1]{fontenc}
\usepackage[margin=1in]{geometry}
\usepackage{amsmath,amssymb,amsthm}
\usepackage{booktabs}
\usepackage{enumitem}
\usepackage{xcolor}
\usepackage[numbers]{natbib}
\usepackage[hidelinks]{hyperref}
\usepackage{float}
\newfloat{algorithm}{t}{loa}
\floatname{algorithm}{Algorithm}

\theoremstyle{definition}
\newtheorem{definition}{Definition}
\newtheorem{proposition}{Proposition}
\newtheorem{remark}{Remark}

\newcommand{\todo}[1]{\textcolor{red}{[\textbf{TODO:} #1]}}
\newcommand{\R}{\mathbb{R}}

\title{\textbf{Morphogenetic Aspect Networks:}\\[2pt]
Self-Evolving Agents via Reward-Modulated Structural Plasticity\\
over an Externally Reasoning LLM}

\author{Linchuan Wang\\
\small HyperdustLabs\\
\small \texttt{linchuan.wang@hyperdustlabs.com}}

\date{Draft --- \today}

\begin{document}
\maketitle

\begin{abstract}
Aspect-Oriented Programming (AOP) was proposed as a biologically inspired
paradigm in which an \emph{aspect} models a neuron, a \emph{pointcut} a
content-conditioned routing primitive, and an \emph{advice} a signal
transduction; a network of aspects (``aspect-of-aspect'') then models a
neural network. When the woven substrate is deterministic code, weaving
reliably changes behavior. The contemporary realization of this idea
weaves into a \emph{stochastic} large language model (LLM) instead, and the
guarantee degrades: most weaving becomes prompt-text injection whose effect
the model may or may not honor. We argue that the missing piece is exactly
the one the original program left open---\emph{synaptic plasticity}---and
that it must be reconceived at the level of \emph{structure}, not weights.
We present \textbf{Morphogenetic Aspect Networks (MAN)}, a formal nucleus in
which an agent is a \emph{stochastic graph-rewriting system}: a typed,
dynamic aspect graph whose fast dynamics perform one inference step
(routing $\to$ weaving $\to$ external-LLM effecting under fixed topology),
and whose slow dynamics \emph{grow and reshape the graph itself}. Crucially,
both are governed by a single variational functional $F=\text{surprise}+\beta\cdot\text{complexity}$:
each structural rewrite fires with rate $\propto\exp(-\Delta F/T)$, so
``learning'' is gradient-free structure search rather than weight descent on
a fixed architecture---a concrete instance of the hypothesis that
\emph{topology and learning are unified}. We specify the state, the credit
assignment that supplies $F$'s error term, and both generative rewrite
primitives (\emph{differentiation/split} and \emph{abstraction/lift}) in full,
and we prove conservation
and behavior-preservation properties that keep self-evolution bounded.
Section~\ref{sec:eval} reports \textbf{Phase~I internal validity} on graded
in-repo fixtures (H1--H5): mechanisms per spec, not application-scale
Self-Evolving proof (Phase~II deferred).
\end{abstract}

%==============================================================
\section{Introduction}
%==============================================================

\paragraph{From separation of concerns to a nervous system.}
The separation-of-concerns principle~\citep{dijkstra1976} and its
realization in Aspect-Oriented Programming~\citep{kiczales1997} modularize
\emph{cross-cutting} behavior that object-oriented decomposition cannot
localize. An earlier line of work observed that biological nervous systems
handle cross-cutting (``non-orthogonal'') concerns more gracefully than
mechanistic paradigms do, and proposed AOP as a candidate
\emph{biologically inspired} programming paradigm, with the mapping
aspect\,$\leftrightarrow$\,neuron, pointcut\,$\leftrightarrow$\,axon,
advice\,$\leftrightarrow$\,signal transduction,
joinpoint\,$\leftrightarrow$\,synapse, object\,$\leftrightarrow$\,effector,
and a three-layer ``evolvable software'' architecture
(organs/spinal-cord/cortex)~\citep{wang2004waosd}. That work flagged one
consideration as unsolved: \emph{synaptic plasticity}, which it argued would
require dynamic AOP.

\paragraph{A note on timing.}
This mapping was articulated in 2004~\citep{wang2004waosd}---over a decade
before content-conditioned routing became central to mainstream architectures
through attention~\citep{vaswani2017} and sparsely-gated
mixture-of-experts~\citep{shazeer2017}, and well before the present wave of
LLM agents. The proposal thus anticipated routing-based, conditionally
computed architectures; what it lacked was a \emph{substrate powerful enough
to serve as the effector}. That effector now exists: a pretrained LLM
(together with its action/control surface---tools, memory writes, message
sends) is the first general, open-domain competence engine an aspect network
can \emph{invoke} rather than hand-program. Crucially, this strength buys
\emph{generality, not determinism}: unlike the deterministic effectors of
2004 (a muscle reliably contracts), the LLM is a \emph{strong but stochastic}
effector. Reconstructing reliable behavior around such an effector is exactly
what MAN exists to do (Section~\ref{sec:discussion})---which is what makes the
revival both possible and necessary.

\paragraph{The reliability gap.}
In the original setting the substrate being woven was deterministic code
(static weaving into bytecode), so ``weaving reliably changes behavior''
held: the joinpoint is a real execution point and the advice is real code.
Contemporary agent systems instead weave into an LLM. The effector is no
longer a deterministic object whose response to a transmitter is fixed
(cf.\ a neuromuscular junction); it is a stochastic generator. Consequently
most advice degrades to text appended to a prompt, whose causal effect is
not guaranteed. Empirically, the only \emph{hard} levers are those that
intercept the host's control flow (refuse a tool call, cancel a message,
veto a sub-agent); the rest is \emph{soft} influence bounded by the model's
instruction-following.

\paragraph{Thesis.}
We take the position that the aspect network should be neither another
trainable artificial neural network (ANN) nor a second LLM, but a
\emph{morphological model of the neuron} in software, committed to two
hypotheses: (H-form) \emph{form and function are unified}---the morphology
of an aspect is its computational role; and (H-topo) \emph{topology and
learning are unified}---unlike an ANN, which separates a fixed differentiable
architecture from gradient-based weight search, a morphogenetic substrate
learns by \emph{reshaping its own structure}. Under these hypotheses,
plasticity is \emph{morphogenesis}: growing, splitting, abstracting, and
pruning structure, scored by a single objective. The aspect is well suited
to be the atomic unit of a neuro-symbolic system because it is, in one
object, both \emph{symbolic} (named, inspectable, composable) and
\emph{sub-symbolic} (graded activation, plastic routing).

\paragraph{Contributions.}
\begin{enumerate}[leftmargin=1.4em,itemsep=2pt]
\item A formal nucleus, \textbf{MAN}
  $M=(N,\Downarrow_\text{fast},\Downarrow_\text{slow},F,\kappa,T(\cdot))$,
  casting an agent as a stochastic graph-rewriting system in which
  \emph{a structural rewrite is an inference step in model space}
  (Section~\ref{sec:model}).
\item A deterministic, replay-reproducible \emph{credit field} $\kappa$
  combining eligibility traces (temporal), responsibility weights
  (structural, with a null-player property), and an advantage baseline,
  obeying credit conservation (Section~\ref{sec:credit}).
\item A single variational selection rule: every rewrite fires with rate
  $\propto\exp(-\Delta F/T)$ where $F=\text{surprise}+\beta\,\text{complexity}$,
  unifying weight and structure updates (Section~\ref{sec:slow}).
\item A full specification of both generative primitives---\emph{differentiation
  (split)} and \emph{abstraction (lift)}---each with guard, rewrite, and
  $\Delta F$ score, plus conservation and behavior-preservation properties
  (Sections~\ref{sec:split}--\ref{sec:props}).
\item A falsifiable evaluation protocol with preregistered hypotheses,
  metrics, baselines, and ablations (Section~\ref{sec:eval}).
\end{enumerate}

\paragraph{Honest scope.}
MAN \emph{orchestrates} an external LLM; reasoning is outsourced to the LLM,
while MAN owns structure, governance, and self-evolution. It does not, by
itself, form new continuous representations. Its single irreducible weakness
(Section~\ref{sec:discussion}) is the causal opacity of soft advice over a
stochastic effector---the same gap that motivates the framework.

%==============================================================
\section{Background and Related Work}
%==============================================================

\paragraph{AOP and process-algebraic semantics.}
AOP~\citep{kiczales1997} adds aspect, pointcut, joinpoint, and advice to a
base paradigm. Its dynamic, mobile communication topology has a natural
home in process calculi: CSP~\citep{hoare1985} and the
$\pi$-calculus~\citep{milner1999}, the latter modeling channel names passed
as messages (\emph{mobility})---the formal substrate for a connection graph
that changes at runtime. Process-algebraic foundations of AOP were studied
by Andrews~\citep{andrews2001}.

\paragraph{Conditional computation and attention.}
A pointcut performs content-conditioned routing, the same primitive as
attention~\citep{vaswani2017} and as the router in sparsely-gated
mixture-of-experts~\citep{shazeer2017}. We use this isomorphism at the level
of \emph{routing topology}; we do not claim the running system is a trained
transformer (Section~\ref{sec:discussion}).

\paragraph{Structural learning and morphogenesis.}
Learning that changes \emph{topology} rather than only weights has a real if
minority lineage: neuroevolution of augmenting topologies
(NEAT)~\citep{stanley2002}, developmental encodings, and biological accounts
of neural morphogenesis and synaptic plasticity~\citep{floreano2001}, with
deep roots in morphogenesis itself~\citep{turing1952}. These approaches
sacrifice the differentiable search space that makes gradient descent
powerful---the central tension we inherit (Section~\ref{sec:discussion}).

\paragraph{Objectives without gradients.}
Our selection rule is a free-energy / minimum-description-length
tradeoff~\citep{rissanen1978,friston2010}; structural moves are accepted by a
Metropolis criterion. Credit assignment uses eligibility traces and the
advantage formulation from reinforcement learning~\citep{suttonbarto2018,
williams1992}, and the Shapley value~\citep{shapley1953} as the principled
(if expensive) responsibility distribution. The fast/slow split mirrors the
System-1/System-2 distinction~\citep{kahneman2011} and society-of-mind
decompositions~\citep{minsky1986}.

%==============================================================
\section{The Morphogenetic Aspect Agent}
\label{sec:model}
%==============================================================

\begin{definition}[Agent as a stochastic graph-rewriting system]
An agent is a tuple
\[
M=\bigl(N,\ \Downarrow_\text{fast},\ \Downarrow_\text{slow},\ F,\ \kappa,\ T(\cdot)\bigr),
\]
where $N$ is a typed dynamic aspect graph (state), $\Downarrow_\text{fast}$
is one inference step under fixed topology, $\Downarrow_\text{slow}$ is the
morphogenetic rewrite relation whose moves fire with rate
$\min\!\bigl(1,\exp(-\Delta F/T)\bigr)$, $F$ is a variational functional,
$\kappa$ is the credit field, and $T(\cdot)$ a temperature schedule.
\end{definition}

The governing claim is that \emph{structure and inference are one process}:
a single application of $\Downarrow_\text{slow}$ is a step of inference in
model space.

\subsection{State $N=(A,S,x)$}
\label{sec:state}
\begin{itemize}[leftmargin=1.4em,itemsep=1pt]
\item $A$: aspect nodes. Each $a$ carries a symbolic identity $\tau(a)$, a
  gain $g(a)$, dendrite/axon ports, and a sliding buffer
  $D_a=\{(\varphi_t,\;a_t,\;r_t)\}$ of context features $\varphi_t\in\R^d$,
  activation score $a_t$, and attributed reward $r_t$.
\item $S$: synapses, directed edges $s=(a_i.\text{axon}\to a_j.\text{dendrite})$
  with weight $w(s)\ge 0$ and eligibility $e(s)$.
\item $x$: the current activation state (``membrane potentials'').
\end{itemize}
\textbf{$\pi$-calculus reading.} An aspect is a process, a synapse a shared
channel name, a pointcut a name-matching predicate; mobility (passing
channel names) is the formal home of structural plasticity.
\textbf{Conserved core.} A subset $A_\text{reflex}\subseteq A$ of
deterministic reflex aspects (inhibitory interneurons) is excluded from
random rewriting; it is the invariant boundary (``brainstem'') that keeps
self-evolution safe.

\subsection{Fast dynamics $\Downarrow_\text{fast}$: one inference step}
\label{sec:fast}
A turn is one $\pi$-reduction at fixed topology:
(1) joinpoint context $\varphi$ (the \emph{query}) matches candidate aspects
via pointcuts, each receiving an activation score $a_i$ (the attention
weight); (2) the weaver composes their advice into an injection under a
budget / top-$k$ (the weighted sum of \emph{values}); (3) the external LLM
(effector) consumes injection plus context and emits output/actions, while
$A_\text{reflex}$ deterministically gates at the effecting boundary
(allow/deny/rewrite). System~1 is the LLM at this step; System~2 is the
aspect graph that schedules it.

%==============================================================
\section{Credit Assignment}
\label{sec:credit}
%==============================================================

A deterministic reward $r_t$ (verifier verdict, effector outcome, task
success) is distributed over aspects and synapses by the product of a
\emph{temporal}, a \emph{structural}, and an \emph{advantage} factor:
\begin{equation}
\kappa(a)\mathrel{+}=(r_t-b)\,e_a(t)\,\rho_a(t),
\qquad
\kappa(s)\mathrel{+}=(r_t-b)\,e_s(t).
\label{eq:credit}
\end{equation}
\begin{itemize}[leftmargin=1.4em,itemsep=1pt]
\item \textbf{Eligibility} $e$ (temporal): $e_a\leftarrow\lambda e_a+\alpha\cdot\text{part}_a$
  (and analogously for edges), handling delayed reward; $\lambda$ is the
  credit horizon (synaptic tagging).
\item \textbf{Responsibility} $\rho$ (structural): tier-1
  influence weighting $\rho_i=a_i\,c_i/\sum_j a_j c_j$ with contribution
  $c_i$ large for \emph{hard}-acting advice and discounted for \emph{soft}
  advice; tier-2 sampled counterfactual / Shapley calibration whose
  null-player property assigns $\approx 0$ credit to irrelevant aspects,
  defeating the ``it fired, therefore it gets credit'' fallacy.
\item \textbf{Baseline} $b$: a context-bucketed running mean, so
  $r_t-b$ is an \emph{advantage} (a reward-prediction error); credit is
  proportional to surprise, not raw reward.
\end{itemize}

\begin{proposition}[Credit conservation]
With $\rho$ normalized ($\sum_i\rho_i=1$) and tier-2 satisfying Shapley
efficiency, the per-turn node credit is conserved:
$\sum_a\kappa_a(t)=r_t-b$. Hence credit is a distribution, not an inflation.
\end{proposition}

Node credit drives $\textsc{split}$/neurogenesis; edge credit drives
$\textsc{connect}$/$\textsc{prune}$/$\textsc{reweight}$. Equation~\eqref{eq:credit}
with $\rho\equiv$ co-activation reduces to reward-modulated Hebbian plasticity
with eligibility traces.

%==============================================================
\section{The Objective $F$ and Slow Dynamics}
\label{sec:slow}
%==============================================================

\begin{definition}[Variational objective]
\[
F \;=\; \underbrace{\textsc{Surprise}}_{\text{task error / verification failure}}
\;+\;\beta\cdot\underbrace{\textsc{Complexity}}_{\text{structural description length}} .
\]
\end{definition}
The same $F$ scores reweighting, edge creation, and node differentiation
alike---there is no separate ``architecture search then weight search.''
This is the formal content of (H-topo).

\paragraph{Rewrite grammar.}
Each rule has the form $\text{LHS pattern}+\text{guard}\Rightarrow\text{RHS}$
and fires with rate $\min(1,\exp(-\Delta F/T))$.
\begin{center}
\begin{tabular}{@{}lll@{}}
\toprule
Primitive & Effect & Driver \\
\midrule
\textsc{connect} (synaptogenesis) & add edge & edge credit $\times$ co-activation \\
\textsc{prune} & remove edge & low weight, cold \\
\textsc{reweight} (LTP/LTD) & change $w$ (degenerate rewrite) & edge credit \\
\textbf{\textsc{split}} (differentiation) & one node $\to$ two specialists & node-credit variance \\
\textsc{merge} (abstraction/fusion) & two near-duplicates $\to$ one & redundancy \\
\textbf{\textsc{lift}} (aspect-of-aspect) & coalition $\to$ higher-order aspect & co-firing coalition \\
\bottomrule
\end{tabular}
\end{center}
\textsc{split} and \textsc{lift} are the \emph{generative core}: the only two
that create new symbolic units (differentiation top-down by specialization;
abstraction bottom-up by coalition). The rest are housekeeping.

\paragraph{Morphogens.}
Local sufficient statistics decide \emph{where} to propose moves
(high advice variance $\to$ propose \textsc{split}; strong coalition
co-firing $\to$ propose \textsc{lift}; high co-activation$\times$credit $\to$
propose \textsc{connect}), compressing the combinatorial space; $\Delta F$
decides \emph{whether} to accept.

\paragraph{LLM-proposed aspects are not self-validating.}
The external LLM may be used as a semantic \emph{proposal mechanism}: from
the current joinpoint, output, verifier trace, and local reward history it can
suggest candidate features or aspects (e.g.\ ``requires citation'',
``missing complexity analysis'', ``unsafe tool boundary''). Such proposals do
\emph{not} enter the dynamic aspect graph merely because the LLM names them.
They first live in a provisional pool with low influence and explicit
provenance. Only reward-bearing evidence can promote them: $\kappa$ estimates
whether the proposed feature explains reward variance or advantage, $F$ scores
the structural cost of materializing it, and $\Downarrow_\text{slow}$ accepts
the move only with the usual Metropolis probability
$\min(1,\exp(-\Delta F/T))$. Thus the LLM proposes morphology, but reward
validates morphology. This boundary prevents MAN from degenerating into
``the LLM writes permanent rules for itself''; symbolic structure becomes part
of $N$ only when it lowers expected surprise enough to pay for its complexity.

\paragraph{The end-to-end loop.}
Algorithm~\ref{alg:man} assembles the state (\S\ref{sec:state}), credit
(\S\ref{sec:credit}), and rewrite grammar above into the self-evolution loop.
The agent is \emph{born} as a zygote---the conserved reflex core
$A_\text{reflex}$ plus a single undifferentiated cortex aspect $a_0$ whose
pointcut matches every joinpoint (lines~1--4). Thereafter $\Downarrow_\text{fast}$
runs each turn (route, weave, effect, attribute credit), and every $P$ turns
$\Downarrow_\text{slow}$ proposes morphogen-localized rewrites and accepts each
with the Metropolis rate $\min(1,\exp(-\Delta F/T))$. Competence is not edited
in by a developer; it is the fixed point that the
$\textsc{split}$/$\textsc{lift}$ dynamics grow---the two generative primitives
are specified next (\S\ref{sec:split},~\S\ref{sec:lift}).

\begin{algorithm}[t]
\caption{Morphogenetic self-evolution: $\Downarrow_\text{fast}$ interleaved with $\Downarrow_\text{slow}$.}
\label{alg:man}
\small
\hrule\vspace{3pt}
\begin{enumerate}[leftmargin=2.9em,itemsep=1.5pt,topsep=2pt,parsep=0pt,label=\textup{\arabic*:}]
\item \textbf{Input:} startup directive $x_0$; external LLM effector \textsc{Llm}; reward \textsc{Reward}; temperature $T(\cdot)$; period $P$.
\item \textbf{Output:} evolved aspect graph $N=(A,S,x)$.
\item \emph{// Genesis: zygote $=$ conserved reflex core $+$ one undifferentiated cortex aspect}
\item $A_\text{reflex}\gets\textsc{SeedReflex}()$ \quad\emph{(deterministic, fail-closed; excluded from $\Downarrow_\text{slow}$)}
\item $a_0\gets\textsc{Extract}(x_0)$,\ \ $\mathrm{pc}(a_0)\gets\top$ \quad\emph{(one intent-aligned cortex aspect; matches every joinpoint)}
\item $N\gets(A_\text{reflex}\cup\{a_0\},\ \varnothing,\ x)$ \quad\emph{(no synapses at genesis)}
\item \textbf{for} each turn $t$ with joinpoint context $\varphi_t$ \textbf{do}
\item \hspace{1.2em}\emph{// $\Downarrow_\text{fast}$: one inference step at fixed topology}
\item \hspace{1.2em}$A_t\gets\{a\in A:\mathrm{pc}(a)\text{ matches }\varphi_t\}$ with activation scores $a_i$
\item \hspace{1.2em}$\mathit{inj}\gets\textsc{Weave}(\text{top-}k\text{ of }A_t\text{ by }a_i)$;\ \ $y_t\gets\textsc{Llm}(\varphi_t,\mathit{inj})$ under $A_\text{reflex}$ gate
\item \hspace{1.2em}$r_t\gets\textsc{Reward}(y_t)$,\ \ $b\gets\textsc{Baseline}(\varphi_t)$ \quad\emph{(advantage $=r_t-b$)}
\item \hspace{1.2em}\emph{// credit assignment, Eq.~\eqref{eq:credit}\, ($\textstyle\sum_a\kappa_a=r_t-b$)}
\item \hspace{1.2em}\textbf{for} $a\in A_t$:\ \ $e_a\gets\lambda e_a+\alpha\,\mathrm{part}_a$;\ \ $\kappa(a)\mathrel{+}=(r_t-b)\,e_a\,\rho_a$;\ \ $D_a\gets D_a\cup\{(\varphi_t,a_t,r_t)\}$
\item \hspace{1.2em}\textbf{for} co-active $s\in S$:\ \ $e_s\gets\lambda e_s+\alpha$;\ \ $\kappa(s)\mathrel{+}=(r_t-b)\,e_s$
\item \hspace{1.2em}\textbf{if} $t\bmod P=0$ \textbf{then} \quad\emph{($\Downarrow_\text{slow}$: morphogenetic rewrite)}
\item \hspace{2.4em}$\textsc{Reweight}(A,S;\kappa)$;\ \ $\textsc{Connect}(S;\kappa)$;\ \ $\textsc{Prune}(S)$ \quad\emph{(housekeeping; no new symbols)}
\item \hspace{2.4em}\emph{// generative: morphogen/LLM proposes} where\emph{; $\Delta F$ decides} whether
\item \hspace{2.4em}$Z_t\gets\textsc{ProposeFeatures}(\varphi_t,y_t,\textsc{Verifier},D)$ \quad\emph{(LLM/heuristic candidates; provisional only)}
\item \hspace{2.4em}\textbf{for} $a\in A\setminus A_\text{reflex}$ with $H(a)\ge\theta_H$ and $n(a)\ge n_\text{min}$ \textbf{do}
\item \hspace{3.6em}$\pi\gets\textsc{BestPartition}(D_a\cup Z_t)$ \quad\emph{(maximize $G(a)$ over observed/proposed features; \S\ref{sec:split})}
\item \hspace{3.6em}\textbf{if} $\textsc{SplitGuard}(a,\pi)\wedge\textsc{Accept}(\Delta F_{\textsc{split}})$ \textbf{then} $\textsc{Split}(a,\pi)$:\ $a\!\to\!a_1,a_2$ with $\mathrm{pc}\wedge\pi,\,\mathrm{pc}\wedge\neg\pi$
\item \hspace{2.4em}\textbf{for} coalition $K\in\textsc{Coalitions}(\text{co-activation graph})$ \textbf{do}
\item \hspace{3.6em}\textbf{if} $\textsc{LiftGuard}(K)\wedge\textsc{Accept}(\Delta F_{\textsc{lift}})$ \textbf{then} $\textsc{Lift}(K)$:\ coalition $\to$ higher-order $a^\star$ (\S\ref{sec:lift})
\item \hspace{2.4em}\textbf{for} near-duplicate $(a_i,a_j)$:\ \textbf{if} $\textsc{Accept}(\Delta F_{\textsc{merge}})$ \textbf{then} $\textsc{Merge}(a_i,a_j)$
\item \hspace{2.4em}consolidate: lift high-credit aspects into $A_\text{reflex}$; archive cold aspects; clear consumed $D_a$
\item \textbf{return} $N$
\item \emph{// \textsc{Accept}$(\Delta F)$: \textbf{true} w.p.\ $\min(1,\exp(-\Delta F/T))$ \,(Metropolis; $T\!\to\!0\Rightarrow$ hard $\Delta F<0$)}
\end{enumerate}
\vspace{1pt}\hrule
\end{algorithm}

%==============================================================
\section{The \textsc{split} Primitive (worked specification)}
\label{sec:split}
%==============================================================

\textsc{split} fires when an aspect is forced to serve two distinguishable
sub-contexts demanding conflicting advice. Over the buffer $D_a$ define the
reward heterogeneity $H(a)=\mathrm{Var}_t[r_t]$ and, for the best
deterministic partition $\pi$ (axis-aligned threshold stumps over each of the
$d$ features plus the leading principal direction; $O(d\,W\log W)$), the
variance-reduction gain
\begin{equation}
G(a)=\mathrm{Var}[r]-\bigl(p_1\mathrm{Var}[r\mid C_1]+p_2\mathrm{Var}[r\mid C_2]\bigr).
\end{equation}

\paragraph{Guard (all required).}
$H(a)\ge\theta_H$; \quad $G(a)/H(a)\ge\theta_\text{sep}$; \quad
$n(a)\ge n_\text{min}$; \quad $|\bar r_1-\bar r_2|\ge\delta$.
The separability ratio rejects splitting on irreducible noise; the sample
and effect-size floors reject splitting on chance.

\paragraph{Rewrite (RHS).}
Replace $a$ by $a_1,a_2$ with
$\mathrm{pc}(a_1)=\mathrm{pc}(a)\wedge\pi$,
$\mathrm{pc}(a_2)=\mathrm{pc}(a)\wedge\neg\pi$; route each incoming synapse to
the child(ren) whose contexts it drove (split weight when ambiguous, then
annealed by \textsc{reweight}/\textsc{prune}); both children inherit outgoing
targets with differentiated gain/advice; children keep named identities
$\tau(a)\!\mid\!C_1,\ \tau(a)\!\mid\!C_2$.

\paragraph{Score.}
$\Delta\textsc{Complexity}=L(\pi)+L(\text{node})+L(\text{synapses})$
(exact). The benefit term has a cheap log-only estimator
$\Delta\textsc{Error}\approx-\eta\,G(a)$ (variance reduction after credit
cleaning, computed by replay), calibrated on the cold path by sampled
counterfactual instantiation. Accept iff
$\Delta F=\Delta\textsc{Error}+\beta\,\Delta\textsc{Complexity}<0$.

%==============================================================
\section{The \textsc{lift} Primitive (worked specification)}
\label{sec:lift}
%==============================================================

\textsc{lift} is the dual of \textsc{split}: where \textsc{split} differentiates
one aspect top-down, \textsc{lift} abstracts a recurring \emph{coalition} of
aspects bottom-up into a higher-order aspect $a^\star$ that cross-cuts
them---literally an ``aspect-of-aspect.'' It is the move that creates
\emph{depth}.

Candidate coalitions $K=\{a_1,\dots,a_k\}$ are not enumerated over all subsets;
they are proposed by cheap community detection / frequent-itemset mining on the
co-activation graph (the morphogen of Section~\ref{sec:slow}). For a candidate
$K$ define, over the joint window, the \emph{cohesion}
$C(K)=P(\text{all }K\text{ active}\mid\text{any }K\text{ active})$ and the
\emph{synergy}
\begin{equation}
S(K)=I(\mathbf{1}_K;\,r)-\sum_{i} I(a_i;\,r)
\quad\bigl(\text{equivalently } R^2_\text{joint}-R^2_\text{additive}\bigr),
\end{equation}
the joint reward-information beyond the additive sum of members. $S(K)>0$ means
the coalition carries an \emph{interaction} the parts do not---something worth
abstracting.

\paragraph{Guard (all required).}
$n(K)\ge n_\text{min}$ (support); \quad $C(K)\ge\theta_\text{coh}$ (they truly
co-fire); \quad $S(K)\ge\theta_\text{syn}$ (genuine synergy); \quad
$2\le|K|\le k_\text{max}$; \quad novelty (no existing higher-order aspect
already covers $K$; otherwise defer to \textsc{merge}).

\paragraph{Rewrite (RHS).}
Create $a^\star$ with a symbolic identity $\tau(a^\star)$ derived from the
coalition members' names $\tau(a_1),\dots,\tau(a_k)$; its pointcut
$\mathrm{pc}(a^\star)$ matches when $K$ co-activates (a threshold/conjunction
over $K$). Wire bottom-up aggregation edges $a_i\!\to\!a^\star$ and top-down
modulation edges $a^\star\!\to\!a_i$, seeded from the measured
co-activation/synergy. Members are \emph{retained} (unlike \textsc{split}, which
replaces): \textsc{lift} adds a layer. \textbf{Identity initialization:}
$a^\star$'s coordinating advice/gain starts at the identity (null modulation),
so $\Downarrow_\text{fast}$ behavior is \emph{unchanged at creation}; $a^\star$
acquires influence only later via \textsc{reweight} gated by $\Delta F$
(cf.\ residual/identity initialization in deep networks).

\paragraph{Score.}
$\Delta\textsc{Complexity}=L(a^\star)+L(\mathrm{pc}^\star)+L(\text{hub edges})-L(\text{prunable lateral edges})$:
abstraction can \emph{reduce} description length by replacing an $O(k^2)$
entangled clique among members with $O(k)$ hub edges through $a^\star$. The
benefit has a cheap log-only estimator $\Delta\textsc{Error}\approx-\eta\,S(K)$
(the synergy a coordinator can capture), calibrated on the cold path by
instantiating $a^\star$ with a learned coordinating advice and replaying sampled
coalition-active turns. Accept iff $\Delta F<0$; given identity initialization,
the subsequent growth of $a^\star$'s influence is itself $\Delta F$-gated.

\begin{remark}[Lift vs.\ Merge]
Both respond to co-occurrence, but to opposite regimes. High \emph{redundancy}
(members duplicate one another; $S(K)\approx0$, high correlation) triggers
\textsc{merge} (collapse to one); high \emph{synergy} (members distinct yet
jointly load-bearing; $S(K)>0$) triggers \textsc{lift} (abstract a coordinator
above them). With \textsc{split}, the three primitives form a complete response
to the three regimes of reward-conditioned co-statistics: separable variance,
synergistic coalition, and redundant duplication.
\end{remark}

%==============================================================
\section{Properties}
\label{sec:props}
%==============================================================

\begin{proposition}[Conservation]
Every \textsc{split} satisfies \emph{domain conservation}
$\mathrm{dom}(a_1)\uplus\mathrm{dom}(a_2)=\mathrm{dom}(a)$ and
\emph{flow conservation} (total in/out weight preserved at creation).
\end{proposition}

\begin{proposition}[Behavior preservation]
\textsc{split} is a \emph{refinement}: since
$\mathrm{pc}(a_1)\vee\mathrm{pc}(a_2)\equiv\mathrm{pc}(a)$ and weight is
conserved, on any context outside the rewritten region the agent's
$\Downarrow_\text{fast}$ behavior is unchanged. \textsc{lift} is
behavior-preserving \emph{at creation} by identity initialization (its
coordinator has null effect until subsequently reweighted under $\Delta F$).
Hence both generative primitives change behavior only where doing so lowers $F$.
\end{proposition}

\begin{proposition}[Replay determinism]
The tier-1 quantities (eligibility, responsibility from logged scores and
hard/soft contribution flags, baseline, $\Delta\textsc{Complexity}$, the
variance-reduction estimator) are fixed functions of the activation/reward
log. Given the log and constants $(\lambda,\alpha,\beta,T(\cdot))$, the
structural trajectory is byte-for-byte reproducible. Tier-2 calibration
introduces stochastic LLM replays and is averaged and seeded.
\end{proposition}

\begin{remark}[Three timescales]
$\Downarrow_\text{fast}$ (per turn) runs inference and accumulates
eligibility; a \emph{warm} loop runs \textsc{reweight}/\textsc{connect}/\textsc{prune};
a \emph{cold} loop runs \textsc{split}/\textsc{lift}/\textsc{merge} plus tier-2
calibration and homeostatic normalization. $T$ is annealed (critical
periods), and freshly created structure is refractory for $\tau_\text{ref}$
turns to prevent split$\leftrightarrow$merge thrashing.
\end{remark}

%==============================================================
\section{Evaluation Plan}
\label{sec:eval}
%==============================================================

We preregister five hypotheses (H1--H5) as \textbf{Phase~I internal validity}:
each mechanism works per spec on graded in-repo fixtures. \textbf{Phase~II}
(application-level Self-Evolving in coding/OpenClaw scenarios) is deferred;
the auxiliary demo-tool table below is illustrative, not an external-validity claim.

\subsection{Hypotheses}
\begin{description}[leftmargin=2.2em,itemsep=2pt]
\item[H1 (Efficiency).] As the network matures, the rate of external-LLM
  calls per solved task decreases (symbolic structure absorbs cognition)
  without loss of task success.
\item[H2 (Useful differentiation).] Each accepted \textsc{split} reduces
  reward variance within its child sub-contexts relative to the parent, and
  improves task reward on those sub-contexts.
\item[H3 (Credit necessity).] Replacing responsibility-weighted credit
  (Eq.~\ref{eq:credit}) with pure participation increases the spurious-split
  rate and degrades final reward.
\item[H4 (Hard$>$soft reliability).] Behavioral changes mediated by hard
  (control-flow) aspects are more reliable and more cleanly creditable than
  those mediated by soft (prompt) aspects; the gap widens with task
  stochasticity.
\item[H5 (Bounded self-evolution).] With the conserved core $A_\text{reflex}$
  and refractory/annealing controls, structural size and reward remain stable
  (no runaway growth or collapse) over long horizons.
\end{description}

\subsection{Phase~II protocol: self-evolution capability (H0)}
\textbf{H0 (primary, deferred).} Through morphogenetic learning alone, with no
code changes, the agent's competence on application scenarios it was \emph{not}
hand-built for rises with experience, and the learned structure transfers to
held-out and cross-domain scenarios---growing toward general competence rather
than overfitting one task. H1--H5 (Phase~I) are the mechanism ablations that
explain \emph{why} H0 should hold; passing Phase~I is necessary but not
sufficient for H0.

\noindent\textbf{Learning curve.} Success/reward versus experience (episodes)
with \emph{no developer edits} mid-run; the signature of self-evolution is a
rising curve produced by $\Downarrow_\text{slow}$ alone.

\noindent\textbf{Transfer / generalization.} Evolve on scenario set $A$;
evaluate frozen structure on held-out members $B$ and on $\geq 1$ cross-domain
set. A small $A\!\to\!B$ gap is the testable surrogate for ``general'': one
substrate and one learning law cover new scenarios \emph{without per-scenario
engineering}.

\noindent\textbf{Decisive baseline.} Beyond the Phase~I baselines, add a
\emph{developer-effort-matched hand-iterated} agent (a human iterates
prompts/tools under a fixed effort budget). The headline H0 result is MAN's
learned curve overtaking the static graph and approaching the hand-iterated
agent at \emph{near-zero developer effort}---the operational form of ``no
longer iterated traditionally.''

\noindent\textbf{Breadth and cost.} Report scenario-family \emph{coverage}
(how many distinct families one substrate$+$law grows into) and
\emph{cost-to-competence} (samples/compute to a target success), since
sparse-reward credit assignment (Section~\ref{sec:discussion}) can make the
curve shallow.

\noindent\emph{Scope.} ``General'' is not claimed literally; it is
operationalized as transfer across a diverse, held-out benchmark plus
cross-domain generalization, with breadth reported. Target environment:
coding / OpenClaw scenarios. \todo{fix scenario benchmark, held-out split,
cross-domain set, and the hand-iteration effort budget.}

\subsection{Metrics (Phase I)}
LLM-calls-per-success; task success / reward; within-child reward variance
ratio (pre/post split); \emph{invalid split rate} on noisy bandit (tier-1
splits vs.\ uniform-$\rho$ no split); reliability gap $\rho_\text{hard}-\rho_\text{soft}$
and outcome-variance gap under injected noise; structural edge-span on
256-row soak; replay-determinism (re-run from log $\Rightarrow$ identical
tier-1 structure). We do \emph{not} treat \textsc{merge} counts as a paper
primitive (grammar table only).

\subsection{Task suites}
Graded in-repo suites (regenerate via
\texttt{scripts/generate\_morphogenetic\_validation\_data.py}):
(i) \texttt{r\_t\_bandit.jsonl} (96 rows, partition
\texttt{zone:alpha}/\texttt{zone:beta}) + \texttt{r\_t\_bandit\_noisy.jsonl};
(ii) verifiable tool guard (\texttt{demo-tool-block}, 5 commands $\times$ 8
epochs); (iii) \texttt{r\_t\_soak\_long.jsonl} (256 rows, 8$\times$ bimodal
replay). External SWE-bench-scale benchmarks remain future work.

\subsection{Baselines and ablations}
Baselines: (a) static aspect graph (no $\Downarrow_\text{slow}$); (b) weight-only
plasticity (\textsc{reweight} alone, ablating \textsc{split}/\textsc{lift});
(c) the underlying LLM with a fixed hand-written prompt; (d) an LLM-only
agent with no aspect graph.
Ablations: responsibility off (H3); eligibility horizon
$\lambda\in\{0,\dots,1\}$; complexity weight $\beta$ sweep; temperature
schedule on/off; tier-2 calibration on/off; conserved core on/off (H5).

\subsection{Phase I results (internal validity, v0.3)}
\noindent\emph{Reproduce: \texttt{bash scripts/run-man-paper-experiments.sh};
primary gate file \texttt{INTERNAL\_VALIDITY.md}; full dump
\texttt{report.json} (\texttt{raw.internal\_validity}). All H1--H5 + F1--F3
gates pass on current fixtures.}

\begin{center}
\begin{tabular}{@{}lll@{}}
\toprule
ID & Claim (abbrev.) & Pass \\
\midrule
H1 & CPS decreases, success held & yes \\
H2 & Split $\downarrow$ variance; bandit lift & yes \\
H3 & $\rho$ spread; tier-1 split on noisy bandit & yes \\
H4 & $\rho_\text{hard}>\rho_\text{soft}$; noise sweep & yes \\
H5 & Soak span bounded; reflex core retained & yes \\
F1--F3 & Replay, $\lambda$, $\beta$ guards & yes \\
\bottomrule
\end{tabular}
\end{center}

\subsection{Auxiliary tables (demo tool model)}

\begin{center}
\begin{tabular}{@{}lcccc@{}}
\toprule
Method & Success $\uparrow$ & LLM calls/succ.\ $\downarrow$ & Reliability gap & Struct.\ stability \\
\midrule
LLM-only & 0.50 & 2.00 & --- & 0.00 \\
Fixed hand prompt & 1.00 & 1.00 & 0.00 & 0.00 \\
Static aspect graph & 1.00 & 1.00 & 0.00 & 0.00 \\
Weight-only plasticity & 1.00 & 1.00 & 0.00 & 0.00 \\
\textbf{MAN (full)} & 1.00 & 1.00 & 0.00 & 0.00 \\
\bottomrule
\end{tabular}
\end{center}

\begin{center}
\begin{tabular}{@{}lccc@{}}
\toprule
Ablation & Success / metric & Spurious-split & Notes \\
\midrule
-- responsibility $\rho$ (H3) & 1.00 & 0.72 vs 0.00 & hard$-$soft spread \\
-- responsibility plasticity & tier-1 split & uniform: no split & noisy + uniform $\rho$ \\
Tier-1 replay & 1.00 & --- & $|{\rm resid}|<10^{-15}$, hash stable \\
H4 hard$>$soft ($\rho$) & gap 0.48 & --- & tier-1 on tied activations \\
Tier-2 LOO & splits 0 vs 0 & --- & $\Delta$ split count \\
-- conserved reflex (H5) & stable & edge span 0 & reflex on: 2 aspects \\
H2 bandit lift & 1.00 & --- & $\Delta\bar r = 0.5$ post-partition \\
\bottomrule
\end{tabular}
\end{center}

\noindent\textbf{H1 longitudinal (10 epochs):} MAN LLM calls/success
$2.0 \rightarrow 1.0$; LLM-only flat at $2.33$; static flat at $1.0$.

\begin{center}
\begin{tabular}{@{}lcc@{}}
\toprule
$\lambda$ & metric & pass \\
\midrule
0.00 & 1.25 & yes \\
0.25 & 1.67 & yes \\
0.50 & 2.50 & yes \\
0.75 & 5.00 & yes \\
0.90 & 12.07 & yes \\
1.00 & 40.00 & yes \\
\bottomrule
\end{tabular}
\end{center}
\begin{center}
\begin{tabular}{@{}lcc@{}}
\toprule
$\beta$ & metric & pass \\
\midrule
0.01 & 1.00 & yes \\
0.02 & 1.00 & yes \\
0.05 & 0.00 & no \\
0.10 & 0.00 & no \\
0.50 & 0.00 & yes \\
1.00 & 0.00 & yes \\
\bottomrule
\end{tabular}
\end{center}
\begin{center}
\begin{tabular}{@{}lcc@{}}
\toprule
noise & metric & pass \\
\midrule
0.00 & 0.00 & no \\
0.10 & 0.03 & yes \\
0.20 & 0.03 & yes \\
0.30 & 0.05 & yes \\
0.40 & -0.02 & no \\
\bottomrule
\end{tabular}
\end{center}

\subsection{Threats to validity}
Reward attribution quality bounds split quality (Section~\ref{sec:discussion});
soft-advice causal estimates are high-variance; feature map $\varphi$ limits
expressible partitions; long-horizon credit is intrinsically hard. The
protocol isolates each via the ablations above.

%==============================================================
\section{Discussion}
\label{sec:discussion}
%==============================================================

\paragraph{The central tension.}
Unifying topology and learning (H-topo) is precisely what removes the
differentiable search space that makes gradient descent powerful. MAN
therefore searches over discrete structure, guided by $\Delta F$ rather than
a gradient. Biology pays for structural learning with evolutionary and
developmental time and massive parallelism; the open question is whether
morphogen-guided proposals plus a cheap $\Delta F$ estimator make it
tractable at useful scale.

\paragraph{The irreducible soft spot.}
The causal effect of a soft advice on a stochastic LLM cannot be known
deterministically; counterfactuals only estimate it, with variance. Hence
\emph{hard}-acting aspects are cleanly creditable while \emph{soft} ones are
statistical. This re-derives, from the credit side, the soft/hard
distinction that motivates the framework, and identifies a concrete lever:
converting soft advice into hard control-flow mediation improves both
reliability and credit cleanliness in one move.

\paragraph{What MAN is and is not.}
MAN is a stochastic graph-rewriting system whose rewrite rates are set by a
variational functional; a readable, governable, growable neuro-symbolic
substrate that \emph{orchestrates} an LLM. It is not a gradient-trained
ANN/transformer, and it does not aim to out-reason a monolithic model on raw
inference; its distinctive value is a transformer one can read, edit, govern,
and grow at the architecture level, fusing the organ/spinal-cord/cortex
layers of the original program into one continuous substrate.

%==============================================================
\section{Conclusion}
%==============================================================

We recast the unsolved ``synaptic plasticity'' of the aspect-oriented
biologically inspired program as \emph{structural} plasticity, and gave a
self-contained formal nucleus---Morphogenetic Aspect Networks---in which an
agent is a stochastic graph-rewriting system whose structure and inference
are one variational process. We specified credit assignment, the variational
selection rule, and both generative primitives (\textsc{split} and
\textsc{lift}) in full, with conservation and behavior-preservation
guarantees, and we preregistered a falsifiable
evaluation protocol. The load-bearing open problems are sharp and named:
low-variance, low-cost credit estimation, and the acquisition of the feature
map $\varphi$.

%==============================================================
\section*{Reproducibility and status}
%==============================================================
This is a draft framework paper. Sections~\ref{sec:model}--\ref{sec:props}
are complete. Section~\ref{sec:eval} Phase~I (internal validity) is filled on
current fixtures; Phase~II application benchmarks remain future work.

%==============================================================
\begin{thebibliography}{99}
%==============================================================
\bibitem{dijkstra1976} E.~W.~Dijkstra. \emph{A Discipline of Programming}. Prentice-Hall, 1976.
\bibitem{kiczales1997} G.~Kiczales, J.~Lamping, A.~Mendhekar, et al. Aspect-Oriented Programming. In \emph{ECOOP}, LNCS 1241, 220--242, 1997.
\bibitem{wang2004waosd} L.~Wang, X.~Tang, L.~Zhang. Aspect-Oriented: a Candidate for the Biologically Inspired Programming Paradigm for Neural Networks and Evolvable Software. In \emph{International Workshop on Aspect-Oriented Software Development (WAOSD'2004)}, Beijing, China, September 2004.
\bibitem{hoare1985} C.~A.~R.~Hoare. \emph{Communicating Sequential Processes}. Prentice-Hall, 1985.
\bibitem{milner1999} R.~Milner. \emph{Communicating and Mobile Systems: the $\pi$-calculus}. Cambridge University Press, 1999.
\bibitem{andrews2001} J.~H.~Andrews. Process-Algebraic Foundations of Aspect-Oriented Programming. In \emph{Reflection}, 187--209, 2001.
\bibitem{vaswani2017} A.~Vaswani, N.~Shazeer, N.~Parmar, et al. Attention Is All You Need. In \emph{NeurIPS}, 2017.
\bibitem{shazeer2017} N.~Shazeer, A.~Mirhoseini, K.~Maziarz, et al. Outrageously Large Neural Networks: The Sparsely-Gated Mixture-of-Experts Layer. In \emph{ICLR}, 2017.
\bibitem{stanley2002} K.~O.~Stanley, R.~Miikkulainen. Evolving Neural Networks through Augmenting Topologies. \emph{Evolutionary Computation}, 10(2):99--127, 2002.
\bibitem{floreano2001} D.~Floreano, J.~Urzelai. Neural Morphogenesis, Synaptic Plasticity, and Evolution. \emph{Theory in Biosciences}, 120:225--240, 2001.
\bibitem{turing1952} A.~M.~Turing. The Chemical Basis of Morphogenesis. \emph{Phil.\ Trans.\ R.\ Soc.\ B}, 237:37--72, 1952.
\bibitem{rissanen1978} J.~Rissanen. Modeling by Shortest Data Description. \emph{Automatica}, 14:465--471, 1978.
\bibitem{friston2010} K.~Friston. The Free-Energy Principle: A Unified Brain Theory? \emph{Nature Reviews Neuroscience}, 11:127--138, 2010.
\bibitem{suttonbarto2018} R.~S.~Sutton, A.~G.~Barto. \emph{Reinforcement Learning: An Introduction}. 2nd ed., MIT Press, 2018.
\bibitem{williams1992} R.~J.~Williams. Simple Statistical Gradient-Following Algorithms for Connectionist Reinforcement Learning. \emph{Machine Learning}, 8:229--256, 1992.
\bibitem{shapley1953} L.~S.~Shapley. A Value for $n$-Person Games. In \emph{Contributions to the Theory of Games II}, 307--317, 1953.
\bibitem{kahneman2011} D.~Kahneman. \emph{Thinking, Fast and Slow}. Farrar, Straus and Giroux, 2011.
\bibitem{minsky1986} M.~Minsky. \emph{The Society of Mind}. Simon \& Schuster, 1986.
\end{thebibliography}

\end{document}
